\documentclass[a4paper, serif, 10pt]{moderncv}

%% moderncv
% style and color
\moderncvstyle{banking}
\moderncvcolor{black}

%% page layout/margins
\usepackage[top=2.5cm, bottom=2.5cm, left=1.5cm, right=1.5cm]{geometry}

\nopagenumbers{}

%% Babel config for Indonesian (Bahasa Indonesia) language
% ref:
% - https://latex3.github.io/babel/guides/locale-indonesian.html
\usepackage[indonesian]{babel}

%% fontspec
\usepackage{fontspec}

\IfFontExistsTF{Linux Libertine}{
  \setmainfont[Ligatures={TeX, Common}, Numbers=OldStyle]{Linux Libertine}
}{}
\IfFontExistsTF{Linux Libertine O}{
  \setmainfont[Ligatures={TeX, Common}, Numbers=OldStyle]{Linux Libertine O}
}{}

%% CTeX
% CJK support, used to show CJK characters in the resume
%
% - fontset=none: disable builtin fontset but instead set the CJK font manually
% - heading=false: disable ctex heading
% - punct=kaiming: use kaiming punctuations styles for CJK
% - scheme=plain: use plain scheme, do not override \normalsize font size
% - space=auto: space settings for CJK characters
%
% ref:
% - http://ctan.mirrorcatalogs.com/language/chinese/ctex/ctex.pdf
\usepackage[UTF8, heading=false, punct=kaiming, scheme=plain, space=auto]{ctex}

\IfFontExistsTF{Noto Serif CJK SC}{
  \setCJKmainfont{Noto Serif CJK SC}
}{}
\IfFontExistsTF{Noto Sans CJK SC}{
  \setCJKsansfont{Noto Sans CJK SC}
}{}

%% line spacing
\usepackage{setspace}
\setstretch{1.125}

%% URL styling - use normal text instead of monospace
\urlstyle{same}

% Auto-underline all links
% Defer until \begin{document} so hyperref is already loaded
\AtBeginDocument{%
  \let\oldhref\href
  \renewcommand{\href}[2]{\oldhref{#1}{\underline{#2}}}%
}

%% Patch \httplink and \httpslink to handle full URLs
% The original moderncv macros always prepend http:// or https:// to URLs.
% This causes issues when the URL already has a protocol (e.g., https://example.com)
% resulting in malformed URLs like https://https://example.com.
% This patch checks if the URL already contains "://" and uses it directly if so.
% Note: We use \str_set:Nx (x-expansion) to expand macros like \@homepage first.
\ExplSyntaxOn
\str_new:N \g_yamlresume_url_str

\RenewDocumentCommand{\httpslink}{O{}m}{
  \str_set:Nx \g_yamlresume_url_str {#2}
  \str_if_in:NnTF \g_yamlresume_url_str {://}
    { \url{#2} }
    { \url{https://#2} }
}

\RenewDocumentCommand{\httplink}{O{}m}{
  \str_set:Nx \g_yamlresume_url_str {#2}
  \str_if_in:NnTF \g_yamlresume_url_str {://}
    { \url{#2} }
    { \url{http://#2} }
}
\ExplSyntaxOff

\name{Dewi Anggraini}{}
\title{Manajer Keberhasilan Pelanggan}
\phone[mobile]{+62 812-3456-7890}
\email{dewi.anggraini@example.com}
\homepage{https://dewanggraini.id}

\address{Jl. Sudirman No. 45, Jakarta, DKI Jakarta, Indonesia, 12190}{}{}

\extrainfo{{\small \faLinkedin}\ \href{https://www.linkedin.com/in/dewianggraini}{@dewianggraini} {} {} {} • {} {} {} 
{\small \faTwitter}\ \href{https://twitter.com/dewianggraini}{@dewianggraini} {} {} {} • {} {} {} 
{\small \faInstagram}\ \href{https://www.instagram.com/dewianggraini}{@dewianggraini}}

\begin{document}

\maketitle

\section{Informasi Dasar}

\cvline{}{\begin{itemize}
\item Manajer Keberhasilan Pelanggan dengan pengalaman 7+ tahun dalam membangun hubungan pelanggan jangka panjang dan meningkatkan retensi.
\item Terampil dalam onboarding pelanggan, analisis data, dan strategi pengurangan churn untuk klien enterprise.
\item Berpengalaman memimpin tim lintas fungsi dan menggunakan CRM seperti Salesforce dan Zendesk.
\item Bersemangat menciptakan pengalaman pelanggan yang positif dan mendorong pertumbuhan pendapatan melalui adopsi produk.
\end{itemize}}
\section{Pendidikan}

\cventry{Sep 2014–Jul 2018}
        {Sarjana, Ilmu Komunikasi, Nilai: 3.75}
        {Universitas Indonesia}
        {\url{https://www.ui.ac.id}}
        {}
        {\begin{itemize}
\item Lulus dengan predikat cum laude dan fokus pada komunikasi bisnis serta perilaku konsumen.
\item Aktif dalam organisasi kemahasiswaan sebagai koordinator hubungan alumni.
\item Menyelesaikan skripsi tentang pengaruh layanan pelanggan terhadap loyalitas.
\end{itemize}
\textbf{Mata Kuliah}: Komunikasi Interpersonal, Manajemen Hubungan Pelanggan, Psikologi Konsumen, Riset Pemasaran, Komunikasi Organisasi, Strategi Komunikasi Digital}
\section{Pengalaman Kerja}

\cventry{Jan 2021–Sekarang}
        {Manajer Keberhasilan Pelanggan}
        {PT Solusi Digital Nusantara}
        {\url{https://www.solusidigital.co.id}}
        {}
        {\begin{itemize}
\item Memimpin tim 8 Customer Success Manager untuk portofolio 120+ klien enterprise dengan nilai ARR Rp 50 miliar.
\item Meningkatkan retensi pelanggan dari 82\% menjadi 94\% dalam 18 bulan melalui program onboarding terstruktur dan business review triwulanan.
\item Mengembangkan playbook keberhasilan pelanggan yang menurunkan waktu onboarding rata-rata dari 45 hari menjadi 21 hari.
\item Berkolaborasi dengan tim produk dan penjualan untuk mengidentifikasi peluang upsell, menghasilkan pertumbuhan pendapatan 27\% YoY.
\item Menganalisis data NPS dan CSAT untuk mengidentifikasi area perbaikan, meningkatkan skor NPS dari 32 menjadi 58.
\end{itemize}
\textbf{Kata Kunci}: Retensi Pelanggan, Onboarding, CRM, Analisis Data, Kepemimpinan Tim}

\cventry{Agu 2018–Des 2020}
        {Spesialis Keberhasilan Pelanggan}
        {PT Karya Inovasi Bangsa}
        {\url{https://www.karyainovasi.co.id}}
        {}
        {\begin{itemize}
\item Mengelola 60+ akun pelanggan UMKM dan menengah, memastikan adopsi produk dan kepuasan pelanggan.
\item Merancang kampanye edukasi pelanggan yang meningkatkan penggunaan fitur utama sebesar 40\%.
\item Menangani eskalasi pelanggan dengan resolusi 95\% dalam waktu kurang dari 24 jam.
\item Membuat laporan kinerja pelanggan bulanan dan merekomendasikan strategi perbaikan.
\end{itemize}
\textbf{Kata Kunci}: Akun Pelanggan, Adopsi Produk, Edukasi, Eskalasi}

\cventry{Jul 2016–Jul 2018}
        {Staf Dukungan Pelanggan}
        {PT Media Cipta Global}
        {\url{https://www.mediacipta.co.id}}
        {}
        {\begin{itemize}
\item Menjawab pertanyaan pelanggan melalui telepon, email, dan obrolan langsung dengan tingkat kepuasan 90\%.
\item Mendokumentasikan masalah pelanggan dan bekerja sama dengan tim teknis untuk penyelesaian.
\item Membantu dalam pelatihan karyawan baru mengenai prosedur dukungan pelanggan.
\end{itemize}
\textbf{Kata Kunci}: Dukungan Pelanggan, Komunikasi, Dokumentasi}
\section{Bahasa}

\cvline{Indonesia}{Bahasa Ibu / Bilingual \hfill \textbf{Kata Kunci}: Bahasa Ibu}
\cvline{Inggris}{Tingkat Profesional Penuh \hfill \textbf{Kata Kunci}: TOEFL 105, IELTS 7.5}
\cvline{Mandarin}{Tingkat Kerja Terbatas \hfill \textbf{Kata Kunci}: HSK 3}
\section{Keahlian}

\cvline{Manajemen Hubungan Pelanggan}{Ahli \hfill \textbf{Kata Kunci}: CRM, Salesforce, Zendesk, HubSpot, Retensi}
\cvline{Analisis Data}{Mahir \hfill \textbf{Kata Kunci}: Excel, SQL, Tableau, Metrik CS, NPS}
\cvline{Komunikasi}{Ahli \hfill \textbf{Kata Kunci}: Presentasi, Negosiasi, Bahasa Indonesia, Bahasa Inggris}
\cvline{Kepemimpinan}{Mahir \hfill \textbf{Kata Kunci}: Manajemen Tim, Mentoring, Kolaborasi Lintas Fungsi}
\section{Penghargaan}

\cventry{Des 2022}
        {PT Solusi Digital Nusantara}
        {Karyawan Terbaik Tahun 2022}
        {}
        {}
        {Diberikan kepada manajer yang menunjukkan kinerja luar biasa dalam meningkatkan retensi pelanggan dan kepemimpinan tim.}

\cventry{Nov 2020}
        {PT Karya Inovasi Bangsa}
        {Penghargaan Inovasi Layanan Pelanggan}
        {}
        {}
        {Mengakui kontribusi dalam merancang program edukasi pelanggan yang meningkatkan adopsi produk secara signifikan.}
\section{Sertifikat}

\cventry{Mar 2022}
        {SuccessCOACHING}
        {Certified Customer Success Manager (CCSM)}
        {\url{https://www.successcoaching.com/certification}}
        {}
        {}

\cventry{Jan 2021}
        {Salesforce}
        {Salesforce Certified Administrator}
        {\url{https://trailhead.salesforce.com/certification}}
        {}
        {}
\section{Publikasi}

\cventry{Jun 2021}
        {Meningkatkan Loyalitas Pelanggan di Era Digital}
        {Jurnal Manajemen Indonesia}
        {\url{https://jurnalmanajemen.id/artikel/loyalitas-digital}}
        {}
        {\begin{itemize}
\item Artikel ini membahas strategi keberhasilan pelanggan untuk mempertahankan loyalitas di pasar digital yang kompetitif.
\item Menekankan pentingnya personalisasi dan komunikasi proaktif.
\end{itemize}}
\section{Referensi}

\cventry{\emaillink[ratna.sari@ui.ac.id]{ratna.sari@ui.ac.id}}
        {Dosen Pembimbing Akademik}
        {Dr. Ratna Sari}
        {+62 811-2233-4455}
        {}
        {Dewi adalah pribadi yang berdedikasi dan memiliki kemampuan komunikasi yang sangat baik. Ia mampu memimpin tim dan menyelesaikan masalah pelanggan dengan efektif.}
\section{Proyek}

\cventry{Feb 2021–Agu 2021}
        {Program onboarding terstruktur untuk klien enterprise}
        {Program Onboarding Pelanggan Enterprise}
        {\url{https://www.solusidigital.co.id/onboarding}}
        {}
        {\begin{itemize}
\item Merancang dan meluncurkan program onboarding 30-60-90 hari yang meningkatkan waktu adopsi.
\item Mengintegrasikan modul pelatihan dan sesi konsultasi.
\item Mencapai skor kepuasan pelanggan 4.8/5.
\end{itemize}
\textbf{Kata Kunci}: Onboarding, Enterprise, Pelatihan}

\cventry{Mar 2022–Nov 2022}
        {Dashboard analitik untuk memantau kesehatan pelanggan secara real-time}
        {Dashboard Kesehatan Pelanggan}
        {\url{https://www.solusidigital.co.id/dashboard}}
        {}
        {\begin{itemize}
\item Mengembangkan dashboard dengan metrik penggunaan produk, tiket dukungan, dan NPS.
\item Membantu tim CS mengidentifikasi risiko churn lebih awal.
\item Meningkatkan efisiensi tim sebesar 30\%.
\end{itemize}
\textbf{Kata Kunci}: Dashboard, Analitik, Churn}
\section{Minat}

\cvline{Membaca}{Buku Bisnis, Psikologi, Pengembangan Diri}
\cvline{Menulis}{Blog, Artikel, Jurnal}
\cvline{Traveling}{Kuliner, Budaya, Alam}
\section{Kegiatan Sukarela}

\cventry{Jan 2022–Sekarang}
        {Mentor Keberhasilan Pelanggan}
        {Komunitas Pelanggan Indonesia}
        {\url{https://www.komunitaspelanggan.id}}
        {}
        {\begin{itemize}
\item Memberikan mentoring kepada profesional muda yang ingin berkarir di bidang keberhasilan pelanggan.
\item Mengadakan workshop tentang strategi retensi dan komunikasi pelanggan.
\item Membantu mengorganisir acara networking tahunan.
\end{itemize}}
    

\cventry{Agu 2019–Des 2020}
        {Relawan Pengajar}
        {Yayasan Pendidikan Anak Bangsa}
        {\url{https://www.yayasanpendidikan.org}}
        {}
        {\begin{itemize}
\item Mengajar keterampilan komunikasi dan literasi digital untuk siswa sekolah menengah.
\item Merancang kurikulum sederhana untuk pengembangan soft skills.
\end{itemize}}
    
\end{document}